\scanpage{74}
1) \uline{Bitriangle} $\Delta$ : \inserted{graphe} isomorphe au \inserted{graphe} produit de deux triangles : $\Delta\simeq s\times t$, \inserted{($\mathrm{Card}\,s=\mathrm{card}\,t=3$)} \inserted{$x=(\alpha,\beta)$ et $x'=(\alpha',\beta')$, \inserted{$x\ne\unsure{y|x'}$} liés ssi $\alpha=\alpha'$ ou $\beta=\beta'$ \deleted{($x\ne x'$)}} \deleted{\unsure{2)}} donc $9$ sommets et $4.9:2=18$ arêtes.
Les triangles $\subset\Delta$ sont les $6$ fibres des deux projections $s\times t\ \substack{\longrightarrow s\\\longrightarrow t}$, \inserted{soit $T$ leur ens,} ils se groupent en deux familles par la relation \textquotedblleft$p\simeq q$'' ssi $p=q$ ou $p\cap q=\varnothing$ ($\Longleftrightarrow\mathrm{card}\,p\cap q\ne1$) L'ens des classes, $\tau$, donne lieu à un rev. d'ordre $3$ $T\longrightarrow\tau$ ($\mathrm{card}\,\tau=2$).
La cat. des $\Delta$ équivaut à celle des ens $T$ de cardinal $\unsure{5|6}$ munis d'une partition de type $(3,3)$, ou encore, des flèches de degré $3$ $T\longrightarrow\tau$ de $\mathrm{Ens}$, avec $\mathrm{card}\,\tau=2$. On associe à $T\longrightarrow\tau$ le graphe $\Delta=\prod_{i\in\tau}T_i$, avec \deleted{les arêtes} $x,y$ \inserted{$\in\Delta$} liés ssi $x\ne y$ et \deleted{\unsure{$x_i$|$y$}} $\exists\,i\in\tau$ tel que $x_i=y_i$ (ou encore, ssi $\exists!\,i\in\tau\quad x_i=y_i$).

2) \uline{Le bloc de triades} C'est un graphe \deleted{\unsure{$\Delta$|$D$}} $D$ tel que le graphe opposé $\overline D$ ($(x,y)$ liés \inserted{dans $\overline D$} ssi $x\ne y$ et $x,y$ \uline{non liés} dans $D$) soit un bitriangle. Ainsi $D=c\times d$, $\mathrm{card}\,c=\mathrm{card}\,\unsure{t|d|T}=\unsure{d|3}$, avec $x=(\lambda,\mu)$ et $\unsure{y'|x'}=(\lambda',\mu')$ liés sss : $\lambda\ne\lambda'$, $\mu\ne\mu'$. Il y a donc.
$\deleted{\frac{\mathrm{card}(D\times D-D)}{2}}\inserted{\binom{\mathrm{card}\,D}{2}}-18=\deleted{\text{arêtes}}18$ arêtes, \unsure{ô|,} tiens tiens, le \unsure{même nb|m. nb} que pour le bitriangle $\overline D$. Il y a $6$ triades i.e. \deleted{\unsure{triplets de}} \inserted{parties} de cardinal $3$ formées de pts mutuellmt non liés, soit $S$ leur ens, se groupant en $2$ familles de $3$, d'où rev. de degré $3$ $S\to\sigma$ ($\mathrm{card}\,\sigma=2$).

\uline{Proposition} Le bloc de triades est isom. à un bitriangle -- i.e. si $\Delta$ \inserted{$=s\times t$} est un bitriangle, $\Delta\simeq\overline\Delta$.

\uline{Dém} On va définir dans $\Delta$ une autre famille $S$ de $6$ parties à trois éléments -- savoir l'ens \inserted{$S$} des triades $\subset\Delta$. Ce sont aussi les parties \inserted{$u$} de $\Delta=s\times t$ \inserted{de cardinal $3$} telles que les deux projections $u\to s$ et $u\to t$ soient injectives (donc bijectives) donc ce sont les graphes des bijections $s\xrightarrow{\sim}t$ -- il y en a en $6$.
Elles se groupent en $2$ familles de trois éléments, en prenant \deleted{les} la relation $\varphi\sim\psi$ ssi $\varphi^{-1}\psi\deleted{^{\unsure{+|\ast}}}$ est une permutation paire de $s$ \deleted{condition} [$S$ \inserted{$=\mathrm{Isom}(s,t)$} est un torseur sous \deleted{$\mathrm{Isom}$ Aut} $\mathrm{Aut}\,s=\unsure{\mathfrak S|\mathscr S}_s$, et on prend les \inserted{$2$} classes sous le sous-groupe d'indice $2$ $\mathfrak S_s^{\unsure{+|\ast}}$; on voit que ce partage ne dépend pas de l'ordre dans lequel on prend $s,t$ \unsure{\ldots|ou}
% archive sheet number: 73
